% !TEX root = ../matan.tex

\section{Числовые ряды. Сходимость ряда, необходимое условие, абсолютная и условная сходимость. Критерий Коши. Гармонический ряд}

\subsection*{Основные определения}

\begin{Definition}{Числовой ряд и его частичные суммы}{}
	Пусть дана последовательность $(a_k)_{k=1}^{\infty}$, $a_k\in\RR$ (или
	$a_k\in\mathbb{C}$). Выражение $\sum_{k=1}^{\infty} a_k$ называется
	\textbf{числовым рядом}, а число
	\[
	S_n:=\sum_{k=1}^{n} a_k
	\]
	--- его $n$-й \textbf{частичной суммой}. Каждому ряду отвечает
	последовательность частичных сумм $(S_n)$.
\end{Definition}

\begin{Definition}{Сходимость ряда}{}
	Ряд $\sum_{k=1}^{\infty} a_k$ называется \textbf{сходящимся}, если
	существует конечный предел его частичных сумм:
	\[
	\exists\,S\in\RR\st \lim_{n\to\infty} S_n=S.
	\]
	Число $S$ называют \textbf{суммой ряда} и пишут $\sum_{k=1}^{\infty}a_k=S$.
	Если конечного предела нет, ряд \textbf{расходится}.
\end{Definition}

\begin{Remark}{}{}
	Исследовать ряд на сходимость --- значит исследовать на сходимость
	\emph{последовательность} его частичных сумм $(S_n)$. Тем самым вся теория
	рядов сводится к теории последовательностей.
\end{Remark}

\subsection*{Необходимое условие сходимости}

\begin{Proposition}{Необходимое условие сходимости}{}
	Если ряд $\sum_{k=1}^{\infty} a_k$ сходится, то $a_k\to0$ при $k\to\infty$.
\end{Proposition}

\begin{proof}
	Пусть $S_n\to S$. Тогда и $S_{n-1}\to S$. Поскольку
	$a_n=S_n-S_{n-1}$, получаем $a_n\to S-S=0$.
\end{proof}

\begin{Remark}{Условие необходимое, но не достаточное}{}
	Из $a_k\to0$ сходимость не следует. Контрпример --- гармонический ряд:
	$\frac1k\to0$, но $\sum\frac1k$ расходится (см.\ ниже).
\end{Remark}

\subsection*{Критерий Коши}

\begin{Theorem}[required]{Критерий Коши для рядов}{}
	Ряд $\sum_{k=1}^{\infty} a_k$ сходится тогда и только тогда, когда
	\[
	\forall\eps>0\ \ \exists N(\eps)\in\NN\ \ \forall n\ge N(\eps)\ \ \forall m\in\NN
	\colon\quad
	\left|\sum_{k=n+1}^{n+m} a_k\right|<\eps.
	\]
\end{Theorem}

\begin{proof}
	Ряд сходится $\iff$ сходится последовательность $(S_n)$. Так как $\RR$
	полно, по критерию Коши для последовательностей это равносильно тому, что
	$(S_n)$ фундаментальна: для любого $\eps>0$ найдётся $N(\eps)$ такое, что
	$|S_p-S_q|<\eps$ при всех $p,q\ge N(\eps)$. Положив $p=n+m$, $q=n$, получаем
	\[
	S_{n+m}-S_n=\sum_{k=n+1}^{n+m} a_k,
	\]
	и условие фундаментальности превращается в указанное неравенство.
\end{proof}

\subsection*{Абсолютная и условная сходимость}

\begin{Definition}{Абсолютная и условная сходимость}{}
	Ряд $\sum_{k=1}^{\infty} a_k$ называется \textbf{абсолютно сходящимся},
	если сходится ряд $\sum_{k=1}^{\infty} |a_k|$. Если сам ряд $\sum a_k$
	сходится, а ряд $\sum |a_k|$ расходится, то $\sum a_k$ называется
	\textbf{условно сходящимся}.
\end{Definition}

\begin{Theorem}{Абсолютная сходимость влечёт сходимость}{}
	Если ряд $\sum_{k=1}^{\infty} |a_k|$ сходится, то и ряд
	$\sum_{k=1}^{\infty} a_k$ сходится.
\end{Theorem}

\begin{proof}
	Применим критерий Коши. Пусть $\sum |a_k|$ сходится; тогда для любого
	$\eps>0$ найдётся $N(\eps)$ такое, что при $n\ge N(\eps)$ и любом $m\in\NN$
	выполнено $\sum_{k=n+1}^{n+m}|a_k|<\eps$. По неравенству треугольника
	\[
	\left|\sum_{k=n+1}^{n+m} a_k\right|
	\le \sum_{k=n+1}^{n+m} |a_k| < \eps.
	\]
	Значит, для $\sum a_k$ выполнен критерий Коши, и ряд сходится.
\end{proof}

\begin{Remark}{}{}
	Обратное неверно: существуют сходящиеся, но не абсолютно сходящиеся
	(условно сходящиеся) ряды, например $\sum_{k=1}^{\infty}\frac{(-1)^{k-1}}{k}$.
\end{Remark}

\subsection*{Гармонический ряд}

\begin{Example}{Расходимость гармонического ряда}{}
	Гармонический ряд $\displaystyle\sum_{k=1}^{\infty}\frac1k$ расходится.

	Сгруппируем члены по блокам, в каждом из которых вдвое больше слагаемых,
	чем в предыдущем:
	\[
	1+\frac12+\underbrace{\left(\frac13+\frac14\right)}_{\ge\,2\cdot\frac14=\frac12}
	+\underbrace{\left(\frac15+\frac16+\frac17+\frac18\right)}_{\ge\,4\cdot\frac18=\frac12}+\cdots
	\]
	В блоке с номером $m$ (то есть в группе из слагаемых
	$\frac{1}{2^{m-1}+1},\dots,\frac{1}{2^m}$) содержится $2^{m-1}$ слагаемых,
	каждое из которых не меньше $\frac{1}{2^m}$, поэтому сумма блока не меньше
	$2^{m-1}\cdot\frac{1}{2^m}=\frac12$. Следовательно,
	\[
	S_{2^m}\ge 1+\frac{m}{2}\xrightarrow[m\to\infty]{}+\infty,
	\]
	то есть частичные суммы неограниченно растут и конечного предела не имеют.
	Ряд расходится.
\end{Example}

